\فصل{نتیجه‌گیری و کارهای آینده}

این پژوهش از یک مسئلهٔ عملی آغاز شد: سامانه‌ای که گفتار فارسیِ نسبتاً پاک را به‌خوبی بازشناسی می‌کند، لزوماً در تماس تلفنی یا \lr{VoIP} قابل اتکا نیست. محدودیت پهنای‌باند، نویز محیطی، اعوجاج کدک، تغییر بهره و افت بسته می‌توانند هم‌زمان رخ دهند و بخشی از نشانه‌های آوایی لازم برای بازشناسی را تغییر دهند یا حذف کنند. از سوی دیگر، دادهٔ برچسب‌خوردهٔ فارسی، چه در شرایط مخابراتی و چه در مقیاس لازم برای پوشش دامنه‌های متنوع، محدود است و گردآوری رونویسی انسانی هزینه و، در مورد تماس‌های واقعی، ملاحظات حریم خصوصی قابل توجهی دارد. راهکار این پایان‌نامه مطالعهٔ دو مسیر مکمل بود: افزایش مقاومت با شبیه‌سازی قابل بازتولید کانال، آموزش چندشرایطی و معماری دوجریانهٔ بهسازی--بازشناسی؛ و افزایش مقیاس داده با گردآوری، قطعه‌بندی، شبه‌برچسب‌گذاری و پالایش صوت بلند فارسی.

دستاورد پژوهش فقط یک مدل نهایی نیست. مجموعه‌ای از اجزای به‌هم‌پیوسته ساخته و ارزیابی شده‌اند: آماده‌سازی چند پیکرهٔ فارسی با قرارداد دادهٔ مشترک، تولید گونه‌های بلند، ساخت جفت‌های پاک--تخریب‌شدهٔ هم‌تراز با فرادادهٔ کامل، تعریف دو خط مبنای کنترل‌شدهٔ ویسپر، ارزیابی فست کانفورمر، پیاده‌سازی مدل همجوشی دوجریانه و ساخت پیکره‌ای ردیابی‌پذیر از کتاب‌های صوتی ایران‌صدا. چهار مسیر اصلی ویسپر-کوچک امکان مقایسهٔ ریزتنظیم معمولی، آموزش چندشرایطی، همجوشی دوجریانه و افزودن دادهٔ شبه‌برچسب‌خورده به همان خط مبنای چندشرایطی را فراهم کردند. آزمایش‌های فست کانفورمر و ویسپر-متوسط نیز مکمل این مقایسه‌ها هستند و به‌ترتیب نتیجهٔ آموزش چندشرایطی در معماری دیگر و عملکرد مدلی با ظرفیت بیشتر را نشان می‌دهند.

\قسمت{جمع‌بندی یافته‌های اصلی}

\زیرقسمت{ارزش آموزش چندشرایطی}

روشن‌ترین یافتهٔ آزمایش‌ها آن است که دادهٔ تخریب‌شده، بدون تغییر معماری ویسپر، مقاومت بازشناسی را به‌طور محسوسی افزایش می‌دهد. دو اجرای تک‌جریانه با بذر، تعداد دوره، نرخ یادگیری، اندازهٔ دسته و روش رمزگشایی یکسان آموزش و دقیقاً روی $21{,}848$ نمونهٔ مشترک ارزیابی شدند. در این مقایسه، افزودن داده‌های تخریب‌شده WER تجمیعی را از $23.27\%$ به $17.83\%$ و CER را از $10.84\%$ به $7.93\%$ کاهش داد. این اعداد به‌ترتیب معادل $23.4\%$ و $26.8\%$ کاهش نسبی خطا هستند.

اهمیت این نتیجه فقط در بهبود عدد تجمیعی نیست. در AGFarsdat، که شامل گفتار تلفنی و \lr{VoIP} واقعی و خارج از فرایند تولید دادهٔ مصنوعی است، WER از $35.63\%$ به $25.80\%$ رسید؛ یعنی $9.83$ واحد درصد کاهش مطلق و حدود $27.6\%$ کاهش نسبی. بنابراین، مدل صرفاً الگوهای شبیه‌ساز را حفظ نکرده و بخشی از مقاومت آموخته‌شده به کانال واقعی منتقل شده است. این مشاهده پاسخ اصلی پژوهش را شکل می‌دهد: وقتی دادهٔ واقعی مخابراتی محدود است، تخریب مصنوعیِ متنوع، کنترل‌شده و قابل ردیابی می‌تواند یک راه عملی برای سازگارکردن بازشناس فارسی با شرایط تلفنی باشد.

این بهبود در همهٔ دامنه‌ها یکنواخت نبود. مدل چندشرایطی در کامن‌وویس پیشرفت بزرگی داشت و در فلورز نیز بهتر شد، اما در پرشین‌اسپیچ و پیکرهٔ گفتار فارسی بهبود نشان نداد. مجموعهٔ نخست فقط $24$ نمونه دارد و برای نتیجه‌گیری قوی کافی نیست؛ افت کوچک مجموعهٔ دوم، در کنار تفاوت سبک و کیفیت داده‌ها، نشان می‌دهد مقاومت کانالی و تعمیم میان‌دامنه‌ای یک مسئله واحد نیستند. افزایش سهم داده‌های تخریب‌شده ممکن است مدل را برای کانال هدف سازگارتر کند، اما تضمین نمی‌کند هر نوع گفتار عمومی یا هر شیوهٔ ضبطی نیز بهتر بازشناسی شود.

نتیجهٔ فست کانفورمر نیز با یافتهٔ اصلی آموزش چندشرایطی هم‌سو بود. افزودن داده‌های تخریب‌شده WER تجمیعی این مدل را از $22.73\%$ به $21.81\%$ کاهش داد و بهبود کوچکی نیز در AGFarsdat ایجاد کرد. این نتیجه نشان می‌دهد سود داده‌های مخابراتی مصنوعی به ویسپر محدود نیست، هرچند ویسپر-کوچک در آزمایش حاضر بهرهٔ بیشتری از این داده‌ها برد.

\زیرقسمت{آنچه مدل همجوشی نشان داد}

مدل دوجریانه برای حل یک تعارض شناخته‌شده طراحی شد: بهساز می‌تواند نویز را کاهش دهد، اما ممکن است هم‌زمان اطلاعات آوایی کم‌انرژی را حذف کند. نگه‌داشتن هر دو نمای اصلی و بهسازی‌شده، رمزگذاری آن‌ها با وزن‌های مشترک، تبادل اطلاعات با توجه متقاطع و انتخاب نرم با یک دروازه باید به مدل اجازه می‌داد در هر ناحیه از نمای مفیدتر استفاده کند. این معماری در فضای لاگ-مل و ویژگی‌های رمزگذار آموزش دید تا هدف بهسازی تا حد امکان با نیاز بازشناس همسو باشد.

نسخهٔ نهایی همجوشی نسبت به نسخهٔ نخست پیشرفت بزرگی داشت و WER را از $24.35\%$ به $18.44\%$ کاهش داد. همچنین نسبت به ویسپر معمولی $4.82$ واحد درصد WER کمتری ثبت کرد و در AGFarsdat با WER برابر $25.38\%$ بهترین عدد مشاهده‌شده را به دست آورد. این نتایج نشان می‌دهند برنامهٔ آموزش اصلاح‌شده و آغاز از یک بازشناس فارسیِ مقاوم، مدل دوجریانه را از یک نمونهٔ اولیهٔ ضعیف به سامانه‌ای رقابتی تبدیل کرده‌اند.

بااین‌حال، هدف سخت‌گیرانه‌تر پژوهش عبور از خط مبنای چندشرایطی بود و این هدف تحقق نیافت. WER تجمیعی همجوشی $0.62$ واحد درصد بیشتر از خط مبنا است و CER آن با مقدار $14.04\%$ به‌وضوح ضعیف‌تر است. مهم‌تر آن‌که میانگین وزن نمای بهسازی‌شده فقط $0.0039$ و وزن نمای اصلی $0.9961$ است. مدل، به‌جای استفادهٔ شرطی از دو منبع، تقریباً همیشه به جریان اصلی بازگشته است. ازاین‌رو، شواهد موجود تأیید نمی‌کنند که بهساز و بلوک همجوشی در شکل فعلی اطلاعاتی پایدار و مکملِ آموزش چندشرایطی فراهم کرده باشند.

این نتیجهٔ منفی از دید پژوهشی آموزنده است. وقتی یک بازشناس قوی از پیش با همان تخریب‌ها آموزش دیده باشد، افزودن یک مسیر بهسازی الزاماً سود بیشتری ایجاد نمی‌کند. ممکن است بازشناس تک‌جریانه پیش‌تر ویژگی‌های مقاوم لازم را آموخته باشد، یا زیان بازسازی بهساز را به سمت شباهت طیفی هدایت کند که برای رمزگشایی متن سودمند نیست. همچنین ممکن است مسیر بهسازی در برخی نواحی مفید باشد، اما دروازه راه‌حل کم‌ریسک‌ترِ اتکا به جریان اصلی را پیدا کند. تشخیص میان این توضیح‌ها به آزمون‌های حذفی نیاز دارد و از وزن میانگین دروازه به‌تنهایی ممکن نیست.

\زیرقسمت{اثر افزایش داده با ایران‌صدا}

مسیر چهارم به‌جای افزودن پیچیدگی معماری، مقیاس و تنوع دادهٔ واقعی فارسی را افزایش داد. از $3{,}777$ کتاب صوتی ایران‌صدا با حدود $6{,}300$ ساعت صوت خام، $250$ ساعت برای پردازش نهایی انتخاب شد. صوت‌ها در سطح منبع ردیابی و اعتبارسنجی، با تشخیص فعالیت گفتاری قطعه‌بندی، با ویسپر-متوسطِ سازگارشده با فارسی شبه‌برچسب‌گذاری، به‌صورت محافظه‌کارانه اصلاح و پالایش، و در سطح کتاب تفکیک شدند. سپس قطعه‌های پذیرفته‌شده بدون حذف هیچ‌یک از اجزای خط مبنا به ترکیب کامل داده‌های معمولی، بلند و تخریب‌شده افزوده شدند. ثابت‌ماندن معماری ویسپر-کوچک و تنظیمات اصلی آموزش، مدل چندشرایطی را به مبنای مستقیم این مقایسه تبدیل می‌کند.

افزودن ایران‌صدا WER تجمیعی را از $17.83\%$ به $17.22\%$ و CER را از $7.93\%$ به $7.70\%$ رساند؛ یعنی به‌ترتیب حدود $3.4\%$ و $2.9\%$ کاهش نسبی خطا. هم‌جهت‌بودن دو معیار نشان می‌دهد دادهٔ افزوده در نتیجهٔ کلی سودمند بوده است، اما این سود از اثر آموزش چندشرایطی کوچک‌تر است و با یک اجرای تک‌بذری نمی‌توان پایداری آماری آن را اثبات کرد.

اثر دادهٔ تازه میان حوزه‌ها یکنواخت نبود. WER نسبت به خط مبنای چندشرایطی در فلورز، پرشین‌اسپیچ و پیکرهٔ گفتار فارسی به‌ترتیب $6.04$، $8.39$ و $5.03$ واحد درصد کاهش یافت و CER نیز در هر سه مجموعه بهتر شد. در مقابل، WER در کامن‌وویس $0.35$ واحد درصد و در AGFarsdat واقعی $0.86$ واحد درصد افزایش یافت. بنابراین، دادهٔ کتاب صوتی بخشی از تعمیم مدل به دامنه‌های عمومی را بهبود داد، اما مقاومت مخابراتی را تقویت نکرد. بااین‌حال، مدل ایران‌صدا در AGFarsdat همچنان $8.97$ واحد درصد WER کمتر از ویسپر معمولی داشت و بخش عمدهٔ سود آموزش چندشرایطی را حفظ کرد. این الگو نشان می‌دهد حجم بیشتر داده به‌تنهایی کافی نیست و همخوانی حوزهٔ دادهٔ افزوده با کاربرد هدف نیز اهمیت دارد.

\زیرقسمت{اثر ظرفیت مدل}

ویسپر-متوسط فارسی با WER برابر $14.36\%$ و CER برابر $6.68\%$ بهترین نتیجهٔ تجمیعی مشاهده‌شده را به دست آورد. این مدل روی همان $21{,}848$ نمونهٔ واجد شرایط سه اجرای ویسپر-کوچک ارزیابی شد، اما ظرفیت آن بیشتر است؛ ازاین‌رو، نتیجه نشان می‌دهد افزایش ظرفیت در تنظیمات حاضر سودمند بوده است، نه این‌که یک روش داده‌ای یا معماری مشخص به‌تنهایی علت بهبود باشد. به همین دلیل، ویسپر-متوسط یک آزمایش تکمیلی برای سنجش ظرفیت است و جای مقایسهٔ کنترل‌شدهٔ چهار مسیر اصلی را نمی‌گیرد.

\قسمت{دستاوردهای پژوهش}

دستاوردهای اصلی پایان‌نامه را می‌توان در چهار سطح خلاصه کرد:

\شروع{فقرات}
\فقره \مهم{داده و بازتولیدپذیری:} خط لوله‌ای مبتنی بر پیکربندی ساخته شد که نویز دیمند، تغییر بهره، محدودیت پهنای‌باند، چرخهٔ واقعی کدک‌های \کد{ffmpeg} و افت بسته را اعمال می‌کند. بذر هر نمونه و همهٔ تصمیم‌های تصادفی در JSONL ثبت می‌شوند. جبران تأخیر کدک، هم‌ترازی طول و هدف پاکِ هم‌پهنای‌باند نیز از آمیخته‌شدن خطای زمانی یا اطلاعات حذف‌شدهٔ کانال با هدف بهسازی جلوگیری می‌کنند.
\فقره \مهم{شاهد تجربی برای ASR فارسی مقاوم:} مقایسهٔ کنترل‌شدهٔ دو ویسپر تک‌جریانه نشان داد آموزش چندشرایطی یک خط مبنای ساده و قوی است. بهبود بزرگ روی AGFarsdat نشان داد شبیه‌سازی کانال می‌تواند به گفتار واقعی مخابراتی تعمیم یابد، هرچند شدت این تعمیم میان مجموعه‌ها متفاوت است.
\فقره \مهم{پیاده‌سازی و ارزیابی صادقانهٔ همجوشی:} یک سامانهٔ دوجریانه با U-Net باقیمانده، مدل‌سازی زمانی در گلوگاه، رمزگذار مشترک، توجه متقاطع دوسویه و دروازهٔ زمان--ویژگی پیاده‌سازی شد. گزارش نتیجه فقط به بهترین سطرها محدود نشد؛ واگرایی WER و CER، فروپاشی دروازه و نرسیدن مدل به خط مبنای قوی نیز به‌عنوان یافته‌های اصلی ثبت شدند.
\فقره \مهم{افزایش مقیاس با نظارت ضعیف:} خط لوله‌ای ردیابی‌پذیر برای گردآوری کتاب‌های صوتی ایران‌صدا، قطعه‌بندی صوت بلند، تولید و اصلاح محافظه‌کارانهٔ شبه‌برچسب، پالایش کیفیت و تفکیک در سطح کتاب ساخته شد. مقایسهٔ کنترل‌شده نشان داد این داده‌ها در مجموع و در سه دامنهٔ عمومی سودمندند، اما اثر آن‌ها حوزه‌وابسته است و به AGFarsdat تلفنی و \lr{VoIP} منتقل نشده است.
\پایان{فقرات}

از کنار هم قرار دادن این چهار سطح، یک نتیجهٔ روش‌شناختی مهم به دست می‌آید: در مسئلهٔ کم‌منبعی مانند ASR فارسی، کیفیت طراحی داده، همخوانی حوزه و انصاف خط مبنا می‌توانند از افزودن اجزای معماری پیچیده اثرگذارتر باشند. معماری جدید باید نه فقط از مدل آموزش‌دیده با دادهٔ معمولی، بلکه از مدلی عبور کند که همان داده‌ها و همان دانش تخریب را بدون سربار معماری دریافت کرده است. به همین ترتیب، ارزش دادهٔ تازه باید با افزودن آن به کامل‌ترین ترکیب پیشین و با گزارش جداگانهٔ حوزه‌ها سنجیده شود؛ زیرا بهبود تجمیعی لزوماً به معنای بهبود در کاربرد هدف نیست.

\قسمت{محدودیت‌های پژوهش}

دامنهٔ نتیجه‌گیری این پایان‌نامه با چند محدودیت داده‌ای، آزمایشی و مدلی مشخص می‌شود.

نخست، شبیه‌ساز حاضر تمام پدیده‌های یک تماس واقعی را پوشش نمی‌دهد. پاسخ ضربه‌ای اتاق، پژواک، نوسان تأخیر، قطع و وصل‌های طولانی، هم‌پوشانی گویندگان و نویز سمت دور مدل نشده‌اند. افت بسته در سطح بسته فقط برای Opus اعمال شده است؛ در سایر کدک‌ها حذف فریم موج یک جانشین تقریبی است و نباید معادل دقیق رفتار شبکه تفسیر شود. همچنین دادهٔ تخریب‌شده عمدتاً از کامن‌وویس ساخته شده و ممکن است تنوع گوینده، سبک گفتار و متن آن نمایندهٔ همهٔ تماس‌های فارسی نباشد.

دوم، هر پیکربندی اصلی با یک بذر گزارش شده و فاصلهٔ اطمینان یا آزمون معناداری برای اختلاف‌ها وجود ندارد. در نتیجه، کاهش بزرگ میان دو مدل تک‌جریانه قانع‌کننده است، اما اختلاف‌های کوچکی مانند $0.42$ واحد درصد در AGFarsdat برای ادعای برتری قطعی کافی نیستند. 

ویسپر و فست کانفورمر از نظر معماری، پیش‌آموزش، رمزگشایی و تنظیمات بهینه‌سازی متفاوت‌اند. بنابراین، اختلاف مطلق آن‌ها را نمی‌توان تنها به نوع معماری نسبت داد و مقایسهٔ قابل اتکاتر برای فست کانفورمر میان دو اجرای معمولی و چندشرایطی همان مدل است.

سوم، معیارها به تفکیک شرایط تخریب گزارش نشده‌اند. بدون تفکیک WER و CER بر اساس SNR، نوع نویز، کدک، پهنای‌باند، نرخ افت بسته و طول گفتار نمی‌توان مشخص کرد کدام بخش خط لوله بیشترین سود را دارد یا مدل همجوشی دقیقاً در چه شرایطی از نمای دوم استفاده می‌کند. نتیجهٔ تجمیعی نیز تحت تأثیر مجموعه‌های بزرگ‌تر است و رفتار دامنه‌های کوچک‌تر را پنهان می‌کند.

چهارم، آزمایش ایران‌صدا فقط $250$ ساعت از حدود $6{,}300$ ساعت صوت گردآوری‌شده را به کار گرفت و کیفیت شبه‌برچسب‌ها با یک نمونهٔ دارای رونویسی انسانی ممیزی نشد. معلم نیز یک مدل بازشناسی است و اصلاح‌گر زبانی به صوت دسترسی ندارد؛ بنابراین، پالایش‌های خودکار خطر انتقال حذف، درج، جانشینی یا اصلاح نادرست را از بین نمی‌برند. اثر این داده تنها با یک اجرای آموزشی سنجیده شد و نبود چند بذر و فاصلهٔ اطمینان برای بهبود تجمیعی کوچک آن اهمیت ویژه دارد. همچنین، داده‌های موجود مشخص نمی‌کنند تفاوت حوزه‌ای نتایج از سبک خواندن، راوی، کیفیت ضبط، توزیع متن، کیفیت شبه‌برچسب یا نسبت اختلاط داده‌ها ناشی شده است.

پنجم، بهساز با معیارهای ادراکی یا سیگنالی مستقل ارزیابی نشده است. این پایان‌نامه عمداً ASR را هدف نهایی قرار داد، اما نبود SI-SDR،‏ PESQ،‏ STOI یا آزمون شنیداری مانع از آن است که دربارهٔ کیفیت صوت بهسازی‌شده داوری شود. افزون بر این، خروجی‌های متنی پرت، به‌ویژه در کامن‌وویس که WER و CER مدل همجوشی رفتاری ناهم‌جهت دارند، هنوز در سطح درج، حذف و جانشینی تحلیل نشده‌اند.

در نهایت، هزینهٔ مدل‌ها اندازه‌گیری نشده است. مدل همجوشی رمزگذار ویسپر را برای دو نما اجرا می‌کند و بهساز و توجه متقاطع نیز سربار دارند. بدون گزارش زمان واقعی، مصرف حافظه و تأخیر، نمی‌توان سنجید که بهبود کوچک آن در بعضی مجموعه‌ها برای یک کاربرد برخط ارزش عملی دارد یا نه.

\قسمت{پیشنهاد برای کارهای آینده}

ادامهٔ پژوهش بهتر است با ارزیابی تشخیصی آغاز شود و سپس به تغییر معماری برسد؛ زیرا در وضعیت فعلی مسئله فقط کمبود ظرفیت مدل نیست، بلکه هنوز معلوم نیست شکست یا موفقیت هر جزء در کدام شرایط رخ می‌دهد.

\زیرقسمت{ارزیابی‌های ضروری کوتاه‌مدت}

\شروع{فقرات}
\فقره خروجی‌های مدل همجوشی باید در چهار حالت «فقط نمای اصلی»، «فقط نمای بهسازی‌شده»، «همجوشی کامل» و «دروازه با وزن ثابت یا تحمیلی» از یک نقطهٔ وارسی واحد ارزیابی شوند. حذف توجه متقاطع و جایگزینی آن با یک دروازهٔ ساده نیز سهم واقعی هر جزء را جدا می‌کند.
\فقره معیارها باید بر اساس نمایهٔ تخریب، SNR، نوع کدک، نرخ و طول افت بسته، پهنای‌باند و طول کلیپ گروه‌بندی شوند. همراه‌کردن این تحلیل با توزیع وزن دروازه روشن می‌کند آیا مدل در شرایط دشوارتر واقعاً به نمای بهسازی‌شده رجوع می‌کند یا نه.
\فقره نمونه‌های دارای بیشترین فاصلهٔ WER و CER باید به‌صورت کیفی بررسی شوند. گزارش تعداد درج، حذف و جانشینی، طول خروجی و چند نمونهٔ ناشناس‌سازی‌شده می‌تواند علت CER بالای همجوشی را آشکار کند.
\فقره برای دادهٔ ایران‌صدا، نمونه‌ای نماینده باید به‌صورت انسانی رونویسی و ممیزی شود تا کیفیت متن خام معلم، اصلاح مدل زبانی و پالایش نهایی جداگانه سنجیده شود. اجرای چندبذری و فاصلهٔ اطمینان نیز باید مشخص کند بهبود تجمیعی محدود این مسیر تا چه اندازه پایدار است.
\پایان{فقرات}

\زیرقسمت{بهبود داده و مدل}

پس از تکمیل ارزیابی‌های بالا، گسترش دادهٔ مخابراتی باید بر نزدیک‌ترکردن شبیه‌ساز به تماس واقعی تمرکز کند. افزودن پاسخ ضربه‌ای اتاق و دستگاه، پژواک، نوسان تأخیر، قطع‌های طولانی‌تر و گفتار هم‌پوشان، همراه با حفظ فرادادهٔ دقیق، پوشش خط لوله را افزایش می‌دهد. استفاده از نویز و پاسخ کانال گردآوری‌شده از محیط‌های فارسی و ساخت دادهٔ تخریب‌شده از چند پیکره، وابستگی به سبک کامن‌وویس را نیز کاهش می‌دهد. در صورت دسترسی قانونی و اخلاقی، یک مجموعهٔ توسعهٔ کوچک از تماس واقعی می‌تواند برای انتخاب مدل به کار رود، درحالی‌که AGFarsdat همچنان به‌عنوان آزمون نهایی دست‌نخورده باقی بماند.

برای صوت بلند، اثر هر مرحله باید با آزمون حذفی جدا شود: آموزش با متن خام معلم در برابر متن اصلاح‌شده، آستانه‌های مختلف پالایش، نسبت‌های متفاوت اختلاط داده و نمونه‌برداری متوازن میان منابع و دامنه‌ها. پس از این ارزیابی، می‌توان بیش از $250$ ساعت کنونی را پردازش کرد و منحنی سود نسبت به حجم داده را سنجید. از آنجا که کتاب صوتی به AGFarsdat منتقل نشد، افزودن صوت بدون برچسبی که از نظر سبک، کانال و تعامل گویندگان به تماس تلفنی و \lr{VoIP} نزدیک‌تر باشد نیز باید به‌عنوان آزمایشی مستقل بررسی شود.

در سطح مدل، هدف بهساز باید بیش از شباهت نقطه‌به‌نقطهٔ لاگ-مل با اطلاعات زبانی همسو شود. تطبیق چندلایهٔ ویژگی‌های رمزگذار، زیان تمایزی برای حفظ واج‌ها، یا آموزش معلم--دانش‌آموز میان نمای پاک و تخریب‌شده گزینه‌های قابل بررسی‌اند. برای جلوگیری از فروپاشی دروازه نیز می‌توان منظم‌ساز آنتروپی، حذف تصادفی یکی از نماها، یا برنامه‌ای برای محدودکردن تدریجی دروازه آزمود. بااین‌حال، هر روش باید در برابر خط مبنای چندشرایطی و یک همجوشی دروازه‌ای ساده سنجیده شود تا سود آن با پیچیدگی غیرضروری اشتباه نشود.

\زیرقسمت{ارزیابی برای استقرار}

اگر هدف نهایی استفاده در سامانهٔ تلفنی برخط باشد، دقت تنها معیار تصمیم نیست. تعداد پارامترهای افزوده، حافظهٔ اوج، زمان پردازش هر ثانیه صوت و تأخیر انتهابه‌انتها باید روی سخت‌افزار هدف اندازه‌گیری شوند. همچنین ارزیابی گفتار پیوسته و طولانی، قطعه‌بندی برخط، تغییر گوینده و بازیابی پس از افت‌های طولانی ضروری است. ممکن است نتیجهٔ چنین ارزیابی‌ای نشان دهد مدل تک‌جریانهٔ چندشرایطی، به‌دلیل دقت بهتر و هزینهٔ کمتر، انتخاب عملی‌تری از همجوشی باشد؛ یا برعکس، یک نسخهٔ سبک‌تر از همجوشی فقط در شرایط بسیار دشوار فعال شود.

\قسمت{سخن پایانی}

نتیجهٔ نهایی این پایان‌نامه آن است که هیچ شکل واحدی از افزایش مقیاس برای همهٔ حوزه‌ها به یک اندازه سودمند نیست. برای مقاوم‌سازی بازشناسی گفتار فارسی در برابر تخریب‌های تلفنی و \lr{VoIP}، داده‌افزایی مخابراتیِ قابل ردیابی و آموزش چندشرایطی از افزودن مستقیم مسیر پیچیدهٔ بهسازی و همجوشی مؤثرتر بود؛ افزودن کتاب صوتی نیز در دامنهٔ مخابراتی سودی ایجاد نکرد. کاهش خطا روی گفتار واقعی AGFarsdat نشان می‌دهد شبیه‌سازی دقیق، حتی با محدودیت‌های موجود، می‌تواند شکاف میان دادهٔ عمومی و کانال واقعی را کاهش دهد.

در سوی دیگر، ایران‌صدا نشان داد صوت واقعیِ شبه‌برچسب‌خورده می‌تواند با حفظ معماری، خطای تجمیعی و خطا در چند دامنهٔ عمومی را کاهش دهد؛ اما سود آن کوچک‌تر و حوزه‌وابسته است. مدل همجوشی نیز یک شکست بی‌حاصل نیست: توسعهٔ نسخهٔ دوم فاصلهٔ آن را با خط مبنای قوی بسیار کم کرد و رفتار دروازه مسئلهٔ مشخصی را برای ادامهٔ پژوهش آشکار ساخت. ویسپر-متوسط بهترین دقت کلی مشاهده‌شده را داشت، ولی مقایسهٔ آن اثر ظرفیت را با سایر عوامل درهم می‌آمیزد. بر پایهٔ مقایسه‌های کنترل‌شده، ویسپر-کوچکِ آموزش‌دیده با ترکیب گفتار معمولی، بلند و تخریب‌شده همچنان روشن‌ترین خط مبنای مقاومت مخابراتی و، در میان مسیرهای کوچک، بهترین توازن میان دقت، سادگی و قابلیت تفسیر را فراهم می‌کند. نتیجهٔ علمی زمانی معتبرتر است که هم پیچیدگی معماری و هم حجم داده تنها در صورت ایجاد سود تکرارپذیر و متناسب با حوزهٔ هدف پذیرفته شوند.
